% STATUS: written 2026-11-05 (Sakai 1976 read in full, Takai 1981 first
% pages, Sakai 1991 preliminaries and contents; excerpts in
% excerpts/Sakai1976OneParameter.txt, excerpts/Takai1981.txt,
% excerpts/Sakai1991Dynamical.txt).
% Additions of 2026-11-06 from Ismert 2018 (pp. 1, 3, 7--8) and 2026-11-07
% from Niknam 1983 (pp. 181--183).
% Sources: Sakai1976OneParameter, Takai1981SakaiProblem, Sakai1991Dynamical,
% Ismert2018Kernel, Niknam1983Closable.
\section{Unbounded derivations and dynamics}\label{sec:banach-operator:unbounded}

Section~\ref{sec:banach-operator:automatic-continuity} closed with
continuity for free: an everywhere defined derivation of a $C^*$-algebra
is bounded and weakly inner. Mathematical physics forces a wider notion:
the derivations of quantum dynamics are defined as infinitesimal
generators of one-parameter subgroups of $*$-automorphisms, and they are
unbounded \cite[p.~427]{Sakai1976OneParameter}. The basic example already
sets the pattern \cite[pp.~427--428]{Sakai1976OneParameter}.

\begin{example}[{\cite[pp.~427--428]{Sakai1976OneParameter}}]\label{ex:banach:ddt}
Let $A=C_0(-\infty,\infty)$ be the $C^*$-algebra of all continuous
functions on $(-\infty,\infty)$ vanishing at infinity. The only bounded
derivation of $A$ is the identically zero mapping. On the other hand, the
differential operator $d/dt$ corresponding to the one-parameter subgroup
of translation operators is an unbounded derivation of $A$. It cannot be
approximated by the zero mapping in any reasonable sense.
\end{example}

The generator behind Example~\ref{ex:banach:ddt} has a general
definition \cite[p.~428]{Sakai1976OneParameter}.

\begin{definition}[{\cite[p.~428]{Sakai1976OneParameter}}]\label{def:banach:generator}
Let $\rho(t)$, $-\infty<t<\infty$, be a strongly continuous one-parameter
subgroup of $*$-automorphisms on a $C^*$-algebra $A$. The strong
derivative of $\rho(t)$ at $t=0$ is a closed, densely defined linear
operator $\delta$ called the \emph{infinitesimal generator} of the
one-parameter subgroup:
$\delta(a)=\lim_{t\to 0}(\rho(t)(a)-a)/t$, and the domain $D(\delta)$
consists of all elements $a$ for which this limit exists.
\end{definition}

An element $a$ is \emph{analytic} if the map $t\mapsto\rho(t)(a)$ is
analytic, and the analytic elements form a dense linear subset of $A$
\cite[p.~428]{Sakai1976OneParameter}. Innerness survives in a restricted
form on the uniformly hyperfinite algebras (Theorem~\ref{thm:banach:uhf})
\cite[Theorem~2, p.~435]{Sakai1976OneParameter}.\begin{theorem}[{\cite[Theorem~2, p.~435]{Sakai1976OneParameter}}]\label{thm:banach:uhf}
Let $A$ be a uniformly hyperfinite $C^*$-algebra and let $\delta$ be a
derivation in $A$ whose domain $D(\delta)$ is a dense $*$-subalgebra of
$A$. Suppose that $D(\delta)$ contains a finite type I-subfactor $F$;
then there is an element $d$ in $A$ such that $\delta(a)=[d,a]$ for all
$a$ in $F$.
\end{theorem}

The whole theory became the book of Sakai on operator algebras in dynamical
systems, devoted to unbounded derivations in $C^*$-algebras
\cite{Sakai1991Dynamical}. The book also fixes the vocabulary used across
this chapter \cite[p.~2]{Sakai1991Dynamical}. A $C^*$-algebra is a \emph{$W^*$-algebra} if it is a dual
space as a Banach space. The $W^*$-algebras are exactly the von
Neumann algebras \cite[p.~2]{Sakai1991Dynamical}.

One question stayed open \cite[p.~202]{Takai1981SakaiProblem}. As a
quantization of spaces, Sakai asked
whether one can find a simple $C^*$-algebra $A$ and a non approximately
bounded pregenerator $\delta_0$ of $A$. Every $*$-derivation
with the same domain should then be a scalar multiple of $\delta_0$ plus an
approximately bounded $*$-derivation \cite[p.~202]{Takai1981SakaiProblem}.
Takai answered a two dimensional form of it on the irrational rotation
algebra $A_\theta$. There are two non approximately bounded
pregenerators $\delta_1,\delta_2$ with the same domain. Each
$*$-derivation of $A_\theta$ with that domain is a combination
$k\delta_1+l\delta_2$ plus an approximately bounded $*$-derivation
\cite[p.~202]{Takai1981SakaiProblem}.

Ismert studies the kernel of an unbounded derivation
\cite[p.~1]{Ismert2018Kernel}. A derivation $\delta$ has kernel
stabilization when $\ker\delta^{n}=\ker\delta$ for every $n$
\cite[p.~1]{Ismert2018Kernel}. She works with the weakly defined derivation
$\delta_{w}$ of $B(H)$ that commutes with a self-adjoint operator $D$
\cite[p.~1]{Ismert2018Kernel}. Its domain consists of the operators that
map the domain of $D$ into itself (Theorem~\ref{thm:banach:kernel-stabilization}) \cite[p.~3]{Ismert2018Kernel}.

\begin{theorem}[{\cite[Theorem~3.5, p.~7]{Ismert2018Kernel}}]\label{thm:banach:kernel-stabilization}
Let $D$ be a self-adjoint operator on a Hilbert space $H$. Then for every
$n\geq 1$ one has $\ker(\delta_{w}^{D})^{n}=\ker\delta_{w}^{D}$.
\end{theorem}

The proof compares the kernel with the commutant of $D$
\cite[p.~7]{Ismert2018Kernel}. Two consequences follow
\cite[p.~7]{Ismert2018Kernel}. The $*$-derivations of Bratteli and Robinson
have kernel stabilization \cite[p.~8]{Ismert2018Kernel}. Also, a pair of
operators that satisfies the Heisenberg commutation relation cannot both be
bounded \cite[p.~8]{Ismert2018Kernel}.

For a simple $C^*$-algebra Niknam connects two of the notions above, the
separating space and the range, through the question of when a derivation can be closed
\cite[p.~181]{Niknam1983Closable}. He takes a derivation $\delta$ defined on a
dense $*$-subalgebra of the algebra, and calls the set of limits of
$\delta(a_n)$ along sequences $a_n\to0$ its separating space. The derivation
is closable exactly when that space is the zero subspace (Theorem~\ref{thm:banach:niknam-closable})
\cite[p.~181]{Niknam1983Closable}.\begin{theorem}[{\cite[Theorem~1, p.~181]{Niknam1983Closable}}]\label{thm:banach:niknam-closable}
Let $A$ be a simple $C^*$-algebra and let $\delta$ be a derivation of $A$. Then
$\delta$ is closable if the range of $\delta$ is not dense in $A$.
\end{theorem}

The proof shows that the separating space is a closed two-sided ideal, hence
either $\{0\}$ or all of $A$, and the density assumption rules out the second
case \cite[p.~182]{Niknam1983Closable}. A $*$-derivation is closable as soon as
one of the two sets $\{a\pm\delta(a)\}$ is not dense
\cite[Corollary~2, p.~182]{Niknam1983Closable}. On a fixed domain, a second $*$-derivation is closed exactly when it is
bounded relative to the first \cite[Theorem~3, p.~182]{Niknam1983Closable}.

The next section returns to bounded
maps on subalgebras of $B(H)$. % claim-ok: link sentence to Section 5, no external fact

For an open connected set in the plane and the Banach algebra of analytic functions on it, the space of bounded derivations is shown not to be reflexive, which is a concrete answer to the reflexivity question for derivation spaces~\cite{WlNonReflexivityOfTheDerivationS2019}.

Kissin sets up a one-to-one correspondence between the unbounded reflexive $*$-derivations of $B(H)$ whose domains are weakly dense and contain the compact operators and pairs $(S,L)$, where $S$ is a symmetric densely defined operator and $L$ is a $J$-orthogonal $*$-reflexive lattice in its deficiency spaces, and analyses the lattices of invariant subspaces of the reflexive algebras these derivations produce~\cite{WlReflexiveStarDerivationsAndLat2006}.

For a closed $*$-derivation $\delta$ of a $C^*$-subalgebra of $B(H)$ that is dominated by a closed symmetric operator $S$ on its domain, Kissin shows that if $S$ has equal finite deficiency indices and the algebra has no finite-dimensional representations, then $S$ has a self-adjoint $\delta$-extension~\cite{WlSymmetricOperatorExtensionsOfU1988}.

Kishimoto and Robinson show that if a closed symmetric derivation $\delta$ of a $C^*$-algebra commutes with a continuous action $\alpha$ of a compact abelian group by $*$-automorphisms and generates a one-parameter group of $*$-automorphisms of the fixed point algebra, then under a density condition on the spectral subspaces $\delta$ generates a one-parameter group of $*$-automorphisms of the algebra~\cite{WlOnUnboundedDerivationsCommutin1982}.
